\documentclass[11pt]{article}
\usepackage{amsmath,amssymb,amsthm,mathtools}
\usepackage[margin=1in]{geometry}
\theoremstyle{plain}
\newtheorem{theorem}{Theorem}
\newtheorem{lemma}{Lemma}
\theoremstyle{remark}
\newtheorem{remark}{Remark}
\newcommand{\E}{\mathbb{E}}\newcommand{\PP}{\mathbb{P}}\newcommand{\R}{\mathbb{R}}
\newcommand{\D}{\mathcal D}\newcommand{\eps}{\varepsilon}

\title{Formalization brief (Skorokhod campaign, 5 of $\sim$6):\\
Aldous's criterion, corrected formulation}
\date{9 July 2026}

\begin{document}
\maketitle

\noindent\textbf{Goal.} Prove Aldous's tightness criterion on
$\D=D([0,1],\R)$, in the OUTER-MEASURE formulation that Task 4's findings
force. Files attached: Basic, Tight, Measurable (and Complete's instances
via import); do not modify them. New library-clean file
\texttt{TypeDDecouplingSkorokhodAldous.lean}. Mathlib imports only; whole
project must build; standard axioms only; \texttt{prop\_aldous} untouched
(next task's wiring).

\section*{Design constraints from Tasks 3--4 (read carefully)}

\begin{itemize}
\item \underline{Never require $w'$ to be measurable.} Task 4 proved the
modulus is NOT captured by rational partitions (irrational-jump
counterexample, documented in the Measurable file) and its level sets are
likely only universally measurable. The assembly must therefore produce, for
each $n$ and $\delta$, an explicitly MEASURABLE event $G_{n,\delta}$
(constructed from the \texttt{crossTime} stopping times and coordinate
evaluations, all measurable by Task 4's core) with
$G_{n,\delta}\subseteq\{w'_{X_n}(\delta)\le2\eps\}$ and
$\PP(G_{n,\delta}^c)$ small. Task 3's bridge
\texttt{isTightMeasureSet\_of\_bdd\_of\_modulus} consumes exactly such outer
bounds (it was engineered on monotonicity/subadditivity) --- check its exact
hypothesis shape and, if it demands slightly different packaging, add an
adapter lemma in the new file rather than re-proving it.
\item \underline{The averaging device is a raw integral, not a measure.}
LeanSearch recon: \texttt{ProbabilityTheory.uniformOn} is the FINITE-SET
uniform (zero on intervals) --- do not touch it. Implement Billingsley
(16.24)ff as
$\frac1{2\delta_0}\int_0^{2\delta_0}\PP(|X(\tau+\delta)-X(\tau)|\ge\eps)\,
d\delta$ with plain (interval) Lebesgue integrals; the swap tools exist:
\texttt{MeasureTheory.intervalIntegral\_integral\_swap} (the exact shape) and
the \texttt{lintegral\_prod}/Tonelli family for the indicator route. No
product-space stopping-time theory is needed anywhere.
\item \underline{No reusable chaining exists} (recon: the
\texttt{IsKolmogorovProcess} internals are not exposed), so the
adjacent-moment criterion is NOT the easy alternative --- full Aldous is
primary; the moment corollary comes from it by Markov.
\end{itemize}

\section*{Tier 1: the two-time estimate (Billingsley's averaging)}

For a $\D$-valued random element $X$ with its coordinate filtration, a
stopping time $\tau\le1$ (Task 3's \texttt{crossTime} instances), and
$\eps,\delta_0>0$, define
$\alpha(\delta_0,\eps):=\sup_{\tau,\ \delta\le\delta_0}
\PP(|X(\tau+\delta)-X(\tau)|\ge\eps)$ (the Aldous quantity; $\tau+\delta$
truncated at $1$). Prove the consecutive-crossing bound: for the iterates
$\tau_{k+1}=\mathrm{crossTime}$ after $\tau_k$ at threshold $\eps$,
\[
  \PP\big(\tau_{k+1}-\tau_k\le\delta_0,\ \tau_{k+1}\le1\big)
  \ \le\ C\,\alpha(2\delta_0,\eps/2)
\]
via the averaged intermediate time: on the event, for EVERY
$\delta\in[\delta_0,2\delta_0]$, $\tau_k+\delta$ lies past $\tau_{k+1}$, and
$|X(\tau_{k+1})-X(\tau_k)|\ge\eps$ forces
$|X(\tau_k+\delta)-X(\tau_k)|\ge\eps/2$ OR
$|X(\tau_k+\delta)-X(\tau_{k+1})|\ge\eps/2$; average the two applications of
the Aldous condition over $\delta$ (Fubini/Tonelli on the indicator ---
follow Billingsley (16.24)--(16.29) closely; both $\tau_k$ and $\tau_{k+1}$
are stopping times, and $\tau_{k+1}$-plus-constant is again admissible in
$\alpha$).

\section*{Tier 2: few crossings $\Rightarrow$ modulus control, on a
measurable event}

From Tier 1 and induction: w.h.p.\ the number of $\eps$-crossings before time
$1$ is at most $K$ and consecutive crossings are $>\delta$-separated ---
an explicitly measurable event $G$ (finitely many conditions on
\texttt{crossTime} values and evaluations). On $G$, the partition by the
crossing times witnesses $w'_X(\delta)\le2\eps$ (path-by-path lemma against
Task 2's \texttt{cadlagModulus}: on each cell between consecutive crossings
the oscillation from the left endpoint is $\le2\eps$ --- mind the every-cell
$>\delta$ convention: merge or extend short terminal cells per the convention
documented in Compact, or bound via a sub-partition argument).

\section*{Tier 3: assembly}

\texttt{aldous\_tightness}: conditions (i) (sup-norm boundedness in
probability) and (ii) ($\alpha(\delta_0,\eps)\to0$ as $\delta_0\to0$,
uniformly in $n$) imply \texttt{IsTightMeasureSet} of the laws --- combine
Tiers 1--2 with the bridge. Then \texttt{aldous\_of\_moment}: if
$\E|X_n(\tau+\delta)-X_n(\tau)|^2\le C\delta$ uniformly, Markov gives (ii).

\begin{remark}[hygiene]
(1) Fallback ladder: full \texttt{aldous\_tightness} $\to$ Tiers 1--2 proved
with the Tier-3 assembly as the sole residual $\to$ Tier 1 alone. No
incorrect statements, no placeholder \texttt{sorry}s in delivered theorems
--- Task 4's discipline.
(2) Adapters over modifications: Basic/Compact/Complete/Tight/Measurable are
frozen.
(3) Report: the exact statement shapes delivered, where the every-cell
convention bit, and any further Billingsley \S16 gaps (the averaging
bookkeeping and the truncation-at-$1$ edge cases are the likely spots).
\end{remark}

\end{document}
